% =====================================================================
%  CHAPTER 3 — PROPOSED SOLUTION (SYSTEM DESIGN AND ALGORITHMS)
%  Graduation thesis: A hybrid (signature + AI) web intrusion detection
%  system operating on Apache HTTP/HTTPS access logs.
%
%  NOTE ON FIGURES / ALGORITHMS / FORMULAS
%  ---------------------------------------
%  To keep this file compilable with only basic LaTeX packages, the flow
%  diagram, the six algorithms, and the formulas are NOT typeset here.
%  Instead, each one is given as a detailed "design specification" block
%  (clearly marked) describing exactly what the figure / pseudocode /
%  equation should contain, so it can be produced by the author or a
%  designer and then dropped in. Search this file for the marker
%  "[[DESIGN SPEC" to find every place that needs an artifact.
%
%  HOW TO USE
%  ----------
%  * Standalone build: compiles on its own (run pdflatex twice for refs).
%    As a standalone document the chapter is numbered 1; inside the thesis
%    it becomes Chapter 3 automatically.
%  * Integration: delete the two "STANDALONE WRAPPER" blocks and keep the
%    content from "\chapter{...}" onward.
% =====================================================================

% --------------------- STANDALONE WRAPPER (begin) --------------------
\documentclass[12pt,a4paper]{report}
\usepackage[utf8]{inputenc}
\usepackage[T1]{fontenc}
\usepackage{lmodern}
\usepackage{amsmath}
\usepackage{booktabs}
\usepackage{array}
\usepackage{enumitem}
\usepackage{graphicx}
\usepackage{url}
\usepackage[hidelinks]{hyperref}
\usepackage{geometry}
\geometry{margin=1in}
\usepackage{textcomp}
\emergencystretch=3em  % absorb minor overfull lines from long identifiers

\newlist{romanlist}{enumerate}{1}
\setlist[romanlist]{label=(\roman*), leftmargin=2.2em, itemsep=2pt}

% A lightweight, dependency-free box used to hold a "design specification"
% (the description of a figure/algorithm/formula to be produced later).
\newenvironment{designspec}[1]
  {\par\medskip\noindent\begin{tabular}{@{}p{0.97\textwidth}@{}}\hline
   \textbf{[[DESIGN SPEC]] #1}\\\hline\end{tabular}\par\nopagebreak\smallskip\itshape}
  {\par\medskip}

\begin{document}
% ---------------------- STANDALONE WRAPPER (end) ---------------------

\chapter{Proposed Solution}
\label{chap:solution}

While Chapter 2 established the problem context and the
theoretical building blocks, this chapter describes the proposed solution
in detail. Section~\ref{sec:overview} gives a bird's-eye view of the entire
processing pipeline, accompanied by a flow diagram, so that the reader can
grasp how a raw log line is transformed into a security verdict. The
subsequent sections then zoom into each stage of the pipeline. First,
Section~\ref{sec:pipeline-feat} describes request parsing and the
extraction of the shared $22$-feature vector. The three detection tiers are
then presented one by one: the signature- and rule-based tier
(Section~\ref{sec:tier1}), the supervised Random Forest tier
(Section~\ref{sec:tier2}), and the unsupervised one-class anomaly tier
(Section~\ref{sec:tier3}). Finally, Section~\ref{sec:consensus} explains the
smart hybrid consensus and the threat-scoring scheme that fuse the three
tiers into a single, explainable decision. For each stage, a step-by-step
description of the underlying procedure complements---but does not
replace---the narrative explanation.

% =====================================================================
\section{Overview}
\label{sec:overview}

The proposed system is a \emph{three-tier hybrid} web intrusion detection
pipeline that consumes Apache access logs and emits, for every request, a
binary verdict (attack or clean), a continuous threat score, and a severity
label. The design principle is to combine three detection paradigms that
have complementary strengths and weaknesses, as motivated in
the related-work survey of Chapter 2: a high-precision signature tier, a
high-accuracy supervised tier, and a high-recall unsupervised tier capable
of flagging previously unseen attacks. A consensus mechanism then reconciles
the three opinions so that the final decision inherits the recall of anomaly
detection while retaining the precision of the supervised model.

Figure~\ref{fig:pipeline} depicts the end-to-end flow of the solution. The
processing of a single log line proceeds through the following main steps,
which correspond directly to the boxes in the diagram:

\begin{romanlist}
  \item \textbf{Parsing and decoding.} Each log line is parsed into its
        constituent fields and its URL is iteratively decoded to undo
        multi-layer encoding.
  \item \textbf{Feature extraction.} The parsed request is turned into a
        $22$-dimensional numeric feature vector shared by the two machine
        learning tiers.
  \item \textbf{Whitelist filtering.} Requests for harmless static
        resources (images, stylesheets, scripts) are recognised and exempted
        from the machine learning tiers to reduce noise.
  \item \textbf{Three-tier detection.} The request is evaluated in parallel
        by Tier~1 (signature and rules), Tier~2 (Random Forest probability),
        and Tier~3 (one-class anomaly model).
  \item \textbf{Consensus and scoring.} The three tier outputs are fused by
        the smart consensus rule into a verdict, and a weighted threat score
        assigns a severity level.
  \item \textbf{Alerting.} The structured result is written to disk in JSON
        Lines format and surfaced on a real-time dashboard.
\end{romanlist}

\begin{figure}[ht]
  \centering
  \fbox{\begin{minipage}{0.92\textwidth}
    \small
    \textbf{Diagram to be produced.} A top-to-bottom flowchart with a
    parallel branch in the middle. The recommended layout is:
    \begin{itemize}[leftmargin=1.4em, itemsep=1pt]
      \item \emph{Top (input):} a parallelogram (data node) labelled
            ``Apache access log (one line per request)''.
      \item A downward arrow to a rectangle ``Parse \& decode URL''.
      \item A downward arrow to a rectangle ``22-feature extraction''.
      \item A downward arrow to a rectangle ``Static-resource whitelist
            filter''.
      \item From this rectangle, \emph{three arrows fan out in parallel} to
            three side-by-side boxes: ``Tier~1: Signature + rules'',
            ``Tier~2: Random Forest ($p_{\text{attack}}$)'', and
            ``Tier~3: One-class LOF''.
      \item The three boxes \emph{converge} with three arrows into a diamond
            (decision node) labelled ``Smart consensus''.
      \item A downward arrow to a rectangle ``Threat score \& severity
            bucket''.
      \item A downward arrow to a parallelogram (output) ``Alert +
            dashboard (JSONL)''.
    \end{itemize}
    Suggested styling: rounded rectangles for processing steps,
    parallelograms for input/output, a diamond for the consensus decision,
    and a light grey fill for the three tier boxes to emphasise the parallel
    stage.
  \end{minipage}}
  \caption{End-to-end processing flow of the proposed three-tier hybrid
           web intrusion detection system.}
  \label{fig:pipeline}
\end{figure}

The overall scanning procedure, implemented in the function
``analyze\_log'', formalises this flow. It iterates over the log lines,
applies the parsing and feature steps, evaluates the three tiers (subject to
the whitelist exemption), and then delegates the final verdict to the
consensus and scoring procedure detailed in Section~\ref{sec:consensus}.
Concretely, the procedure carries out the following steps for the input log
file:

\begin{enumerate}[leftmargin=1.8em, itemsep=2pt]
  \item For each log line, parse the line into its fields and produce the
        decoded URL.
  \item Extract the $22$-dimensional feature vector from the parsed request.
  \item If the request is whitelisted (a static asset, or a request with
        status code $408$), skip the two learning tiers and treat their
        outputs as ``not fired''; otherwise continue.
  \item Evaluate Tier~1 to obtain the set of fired rule identifiers.
  \item Evaluate Tier~2 to obtain the Random Forest attack probability $p$.
  \item Evaluate Tier~3 to obtain the one-class anomaly flag $a$.
  \item Call the consensus procedure with the rule set, $p$, and $a$ to
        obtain the verdict, the threat score, and the severity level.
  \item Write the structured record (verdict, score, severity, contributing
        tiers) to the output file.
\end{enumerate}

In addition to the per-request analysis, the system performs a lightweight
behavioural pass after the loop. It tracks, for every source IP address, the
set of distinct paths requested; an IP that touches at least ten distinct
paths is flagged as a probable directory or port scanner. This complements
the content-based tiers with a simple traffic-pattern heuristic. The same
pipeline is reused unchanged in the real-time monitoring mode, where log
lines are read from standard input and alerts are streamed to the dashboard
as they occur.

% =====================================================================
\section{Request Parsing and Feature Extraction}
\label{sec:pipeline-feat}

\subsection{Parsing and multi-layer URL decoding}

The first stage converts a raw Apache log line into a normalised structure
used by every downstream component. The function ``parse\_log\_entry''
relies on a log parser configured for the Combined Log Format, with a
fallback to the Common Log Format for lines that lack the
``Referer''/``User-Agent'' headers. From the request line it
extracts the HTTP method, the raw URL, and the protocol; it then URL-decodes
the URL and splits it into a path, a query string, and a dictionary of
parameters. The status code, the ``User-Agent'', and the
``Referer'' headers are also retained, since they feed several
features described below. The result is a record with the fields
``path'', ``query'', ``method'', ``status'',
``user\_agent'', and ``referer'', which constitutes the common
input contract for all three tiers.

A subtle but important detail is the handling of encoded payloads. As
discussed in Chapter 2, attackers hide malicious content
behind several layers of URL encoding. The helper ``multi\_decode''
therefore decodes the URL repeatedly, up to three iterations, stopping early
when a further decoding leaves the string unchanged. This recovers the
original payload and, as a by-product, yields the \emph{encoding depth}
(the number of layers that were actually peeled off), which is itself an
informative feature.

\subsection{The shared 22-feature vector}

Both machine learning tiers operate on the same $22$-dimensional feature
vector, produced by the ``extract\_features'' method. Sharing a single
feature space guarantees that the performance gap between the supervised and
unsupervised tiers reflects the ``labelled versus unlabelled'' distinction
alone, rather than any difference in representation. The twenty-two features
fall into three conceptual groups---structural, randomness-and-encoding, and
semantic-and-contextual---which are described feature by feature below; the
corresponding source code is shown in Figure~\ref{fig:code-features}.

\paragraph{Structural group (nine features).} This group captures the coarse
size and shape of the request. It comprises the total URL length
(``len\_url''), counted over the concatenated path and query; the path
depth (``path\_depth''), the number of ``/'' separators in the
path, which grows for deeply nested or traversal-style endpoints; a binary
POST flag (``is\_post'') recording the request method; the parameter
count (``param\_count''), taken as the number of ``='' signs in
the query; the raw count of dangerous metacharacters (such as the
semicolon, pipe, dollar sign, back-tick, parentheses, angle brackets,
quotes, braces, and backslash) in the decoded URL
(``raw\_risk\_count'') together with their density relative to the URL
length (``risk\_ratio''); the number of SQL, XSS, and LFI keywords found
in the decoded URL (``keyword\_hits''); the percent-encoding density,
i.e. the proportion of ``\%'' characters
(``encoding\_ratio''); and the Shannon entropy of the decoded URL
(``entropy''). For a string $s$, the Shannon entropy is
\begin{equation}
  H(s) = -\sum_{c \in \Sigma} \frac{n_c}{|s|}\,\log_2 \frac{n_c}{|s|},
  \label{eq:entropy}
\end{equation}
where $\Sigma$ is the set of distinct characters occurring in $s$, $n_c$ is
the number of occurrences of character $c$, and $|s|$ is the total length of
$s$, so that $n_c/|s|$ is the relative frequency of $c$. A payload that has
been encoded or randomised spreads its characters more uniformly and
therefore attains a higher entropy, which makes Equation~\eqref{eq:entropy} a
compact, signature-independent signal of obfuscation. Taken together, an
over-long URL, a high risky-character density, or an abnormally high entropy
are all hallmarks of an injection or buffer-overflow attempt.

\paragraph{Randomness-and-encoding group (five features).} This group
quantifies obfuscation and the multi-layer encoding evasion discussed in
Chapter 2. It contains the maximum Shannon entropy taken over all individual
parameter values (``param\_value\_max\_entropy''), which isolates a
single high-entropy field even when the rest of the URL looks ordinary; the
length of the longest parameter value (``longest\_param\_value''), which
flags an oversized payload smuggled into one field; the encoding depth
(``encoding\_depth''), the number of URL-decoding passes that were
actually required (from $0$ up to $3$) before the payload stopped changing,
which directly measures how many layers of encoding the attacker stacked; the
consecutive-encoding intensity (``consecutive\_encoding\_intensity''),
the fraction of the URL occupied by runs of consecutive ``\%xx''
sequences, which is high for heavily encoded payloads; and the number of
parameter names drawn from a curated suspicious set such as ``cmd'',
``exec'', ``file'', or ``passwd''
(``param\_names\_anomaly'').

\paragraph{Semantic-and-contextual group (eight features).} This group
encodes domain knowledge about how attacks differ from legitimate traffic. It
comprises the suspicious-User-Agent score (``suspicious\_ua\_score''), a
value in $[0,1]$ that is maximal for a known scanner string (``sqlmap'',
``nikto'', and the like), moderate for an automated client, and high for
a missing User-Agent; the Referer-suspicion score
(``referer\_suspicion''), likewise in $[0,1]$, which rises when the
Referer header itself contains attack tokens or is heavily percent-encoded;
the status-code indicator (``status\_code\_indicator''), a small bucket
of the HTTP response status ($0$ for 2xx, $1$ for 3xx, $2$ for 4xx, $3$ for
5xx) that captures error-inducing requests; the signature score
(``signature\_score''), the number of Tier-1 attack regular expressions
(out of four) that match the decoded payload, which feeds the signature
evidence back into the learning tiers; the risk-keyword density
(``risk\_keyword\_density''), the fraction of word tokens that are
dangerous keywords; the method-and-path combination risk
(``method\_path\_combo\_risk''), which rises for sensitive paths such as
``/admin'' or ``/login'' and for write methods such as ``PUT''
or ``DELETE''; the count of special characters in the path itself
(``path\_special\_chars''), which exposes traversal or injection embedded
in the path rather than in the query; and the unknown-parameter ratio
(``unknown\_param\_ratio''), the fraction of the request's parameter
names that were never observed on that exact endpoint during training---the
feature that exposes the parameter-name-mutation attacks pervasive in the
CSIC dataset.

The extraction procedure first decodes the URL to obtain a canonical payload
and the encoding depth, then computes the three groups of features and
concatenates them into a single vector. The procedure is deterministic and
stateless except for the per-path vocabulary, which is learned once at
training time and consulted by the unknown-parameter-ratio feature. Its
steps are:

\begin{enumerate}[leftmargin=1.8em, itemsep=2pt]
  \item Concatenate the path and query string into the full URL.
  \item Decode the full URL up to three times to recover the payload and to
        record the encoding depth.
  \item Compute the structural features (length, path depth, POST flag,
        parameter count, risky-character ratio, and so on).
  \item Compute the randomness-and-encoding features (decoded-URL entropy,
        maximum parameter-value entropy, encoding depth, and
        consecutive-encoding intensity).
  \item Compute the semantic-and-contextual features (signature score,
        risk-keyword density, header-suspicion scores, and the
        unknown-parameter ratio against the learned vocabulary).
  \item Concatenate the three groups into the final $22$-dimensional vector
        and return it.
\end{enumerate}

\begin{figure}[ht]
  \centering
  \fbox{\parbox{0.9\textwidth}{\centering\vspace{2.2cm}
    \textit{[Insert a screenshot of the ``extract\_features()'' method
    and the ``FEATURE\_NAMES'' list from ``models/ai\_detector.py''
    here.]}\vspace{2.2cm}}}
  % Replace the box above with: \includegraphics[width=0.9\textwidth]{code_extract_features.png}
  \caption{Source code of the 22-feature extraction routine
           (``extract\_features'' in ``ai\_detector.py'').}
  \label{fig:code-features}
\end{figure}

% =====================================================================
\section{Tier 1 --- Signature- and Rule-Based Detection}
\label{sec:tier1}

The first tier implements the misuse-detection paradigm. It is always
active, requires no training data, and is the most trusted signal in the
system because of its very high precision. Conceptually, the tier applies a
battery of pattern matchers to the decoded payload and reports which, if
any, attack rule fired. The matchers cover the attack classes enumerated in
Section~2.4 of Chapter 2: SQL injection, cross-site scripting, path
traversal and local file inclusion, command injection, remote file
inclusion, sensitive-file access, scanner user agents, HTTP method
tampering, and parameter tampering.

\subsection{The rule set}

Each attack class is encoded as a pre-compiled, case-insensitive regular
expression (or, for the two stateful rules, as a small procedure). The nine
rules are summarised in Table~\ref{tab:rules}; their exact definitions in the
source code are shown in Figure~\ref{fig:code-regex}.

\begin{table}[ht]
  \centering
  \caption{The nine Tier~1 detection rules, their matching mechanism, and the
           attack class each targets.}
  \label{tab:rules}
  \renewcommand{\arraystretch}{1.3}
  \setlength{\tabcolsep}{5pt}
  \footnotesize
  \begin{tabular}{|p{3.2cm}|p{6.6cm}|p{2.6cm}|}
    \hline
    \textbf{Rule (identifier)} & \textbf{Matching mechanism (representative patterns)} & \textbf{Target attack} \\
    \hline
    \texttt{SQL\_INJECTION} & \texttt{union select}, \texttt{select \dots from}, \texttt{insert into}, \texttt{drop table}, quoted tautologies such as \texttt{'OR'a'='a}, comment markers \texttt{--}/\texttt{\#}/\texttt{/*}, \texttt{xp\_cmdshell} & SQL injection \\
    \hline
    \texttt{XSS} & \texttt{<script}, \texttt{javascript:}, \texttt{onerror=}, \texttt{onload=}, \texttt{eval(}, \texttt{<img src}, \texttt{alert(} & Cross-site scripting \\
    \hline
    \texttt{LFI\_PATH\_} \texttt{TRAVERSAL} & \texttt{../}, \texttt{..\textbackslash}, \texttt{/etc/passwd}, \texttt{/proc/self}, \texttt{C:\textbackslash Windows}, \texttt{win.ini} & Path traversal / LFI \\
    \hline
    \texttt{COMMAND\_} \texttt{INJECTION} & \texttt{;|\&\&\$>} before \texttt{rm}/\texttt{cat}/\texttt{wget}/\texttt{curl}/\texttt{bash}/\texttt{nc}/\dots, or back-tick command substitution & OS command injection \\
    \hline
    \texttt{RFI} & A parameter value beginning with \texttt{http://}, \texttt{https://}, or \texttt{ftp://} & Remote file inclusion \\
    \hline
    \texttt{SENSITIVE\_} \texttt{FILE\_ACCESS} & \texttt{.env}, \texttt{.git/config}, \texttt{wp-config.php}, \texttt{.bak}, \texttt{.sql}, \texttt{/cgi-bin/}, \texttt{phpinfo.php} & Config/backup file access \\
    \hline
    \texttt{SCANNER\_TOOL} & User-Agent matching \texttt{sqlmap}, \texttt{nikto}, \texttt{nmap}, \texttt{nessus}, \texttt{dirbuster}, \texttt{gobuster}, \texttt{acunetix}, \texttt{hydra} & Automated scanners \\
    \hline
    \texttt{HTTP\_METHOD\_} \texttt{TAMPERING} & An unusual method (\texttt{PUT}, \texttt{DELETE}, \texttt{PATCH}, \texttt{TRACE}, \texttt{CONNECT}) on a non-API path & Method tampering \\
    \hline
    \texttt{PARAM\_} \texttt{TAMPERING} & A parameter name absent from the per-path vocabulary learned during training & Parameter-name mutation \\
    \hline
  \end{tabular}
\end{table}

The first seven rules are purely pattern-based and stateless. Their design
reflects a deliberate effort to balance coverage against false positives. For
example, the SQL-injection expression was broadened beyond the classic
``digit-equals-digit'' tautology (``1=1'') to also match the
\emph{quoted, letter-based} tautologies such as ``'OR'a'='a'' that
dominate the CSIC dataset, because the narrower pattern missed the majority
of real CSIC payloads. The command-injection expression requires a
shell metacharacter (a semicolon, pipe, double ampersand, dollar sign, or
back-tick) \emph{immediately before} a known binary name, which keeps the
benign word ``cat'' in a product description from triggering the rule. The
remote-file-inclusion rule inspects parameter \emph{values} rather than the
whole URL, so that a legitimate path containing ``http'' does not raise
an alert.

The last two rules are stateful: they consult a \emph{per-path parameter
vocabulary} that the system learns at training time, recording for every
endpoint the set of parameter names that ever appeared on it. The
method-tampering rule fires when one of the methods ``PUT'',
``DELETE'', ``PATCH'', ``TRACE'', or ``CONNECT'' is used
on an endpoint that is not part of an API surface, since such methods are
rarely legitimate on ordinary pages. The parameter-tampering rule fires when
a request carries a parameter name that was never observed on that exact path
during training; this is the mechanism that catches the frequent
name-mutation attacks in CSIC, for example renaming ``modo'' to
``modoA'' or ``precio'' to ``precioA''.

\begin{figure}[ht]
  \centering
  \fbox{\parbox{0.9\textwidth}{\centering\vspace{2.2cm}
    \textit{[Insert a screenshot of the regular-expression definitions
    (``SQLI\_RE'', ``XSS\_RE'', ``LFI\_RE'',
    ``CMD\_INJ\_RE'', ``SENSITIVE\_FILE\_RE'',
    ``SCANNER\_UA\_RE'') from ``apache\_log.py'' here.]}\vspace{2.2cm}}}
  % Replace the box above with: \includegraphics[width=0.9\textwidth]{code_regex_rules.png}
  \caption{Source code of the Tier~1 signature regular expressions in
           ``apache\_log.py''.}
  \label{fig:code-regex}
\end{figure}

\subsection{Rule evaluation}

The crucial first step of the procedure is defensive decoding: the payload
is decoded up to three times \emph{before} any regular expression is
applied, so that multi-layer encoded attacks such as a triple-encoded SQL
tautology are normalised back to their canonical form and can be matched.
Two of the rules are stateful and depend on the per-path vocabulary learned
during training. The method-tampering rule fires when an unusual HTTP method
(``PUT'', ``DELETE'', ``PATCH'', ``TRACE'',
``CONNECT'') is used on a non-API endpoint. The parameter-tampering
rule fires when a request carries a parameter name that was never seen on
that exact path during training, which is the mechanism that catches CSIC's
frequent name-mutation attacks (for example
``modo''~$\rightarrow$~``modoA''). The rule-evaluation procedure
proceeds as follows:

\begin{enumerate}[leftmargin=1.8em, itemsep=2pt]
  \item Initialise an empty set of fired rule identifiers.
  \item Form the candidate string by joining the path and the query string.
  \item Decode the candidate string up to three times, stopping early if a
        decoding leaves it unchanged.
  \item Apply each signature matcher (SQL injection, XSS, LFI/path
        traversal, command injection, remote file inclusion, sensitive-file
        access, scanner user agent) to the decoded string; for every match,
        add the corresponding identifier to the set.
  \item Apply the method-tampering rule; if it fires, add its identifier.
  \item Apply the parameter-tampering rule against the per-path vocabulary;
        if it fires, add its identifier.
  \item Return the set of fired rule identifiers.
\end{enumerate}

Figure~\ref{fig:code-detect} shows the implementation of this procedure in
the ``detect\_rule\_based'' function.

\begin{figure}[ht]
  \centering
  \fbox{\parbox{0.9\textwidth}{\centering\vspace{2.0cm}
    \textit{[Insert a screenshot of the ``detect\_rule\_based()''
    function (and the ``multi\_decode'' loop) from
    ``apache\_log.py'' here.]}\vspace{2.0cm}}}
  % Replace the box above with: \includegraphics[width=0.9\textwidth]{code_detect_rule_based.png}
  \caption{Source code of the Tier~1 rule-evaluation routine
           (``detect\_rule\_based'' in ``apache\_log.py'').}
  \label{fig:code-detect}
\end{figure}

As established experimentally in Chapter 2, this tier
attains a precision of $99.61\%$ but a recall of only $24.07\%$. In other
words, when it fires it is almost always correct, but it stays silent on the
majority of attacks. This profile makes it ideal as a high-confidence
trigger within the consensus rule, but unsuitable as a standalone detector,
which is the very reason the two learning tiers are introduced next.

% =====================================================================
\section{Tier 2 --- Supervised Detection with Random Forest}
\label{sec:tier2}

The second tier is a supervised classifier that learns attack patterns from
labelled data. As justified in Chapter 2, the Random Forest is
chosen as the primary model because of its robustness, its tolerance of
unscaled heterogeneous features, and its ability to output a calibrated
attack probability. The classifier is configured with $200$ trees and
balanced class weights and is trained on the labelled portion of the dataset
using the same $22$-feature vector as the other tiers.

Training comprises three conceptual steps. First, the per-path parameter
vocabulary is built from the training logs so that the
unknown-parameter-ratio feature is meaningful. Second, every labelled
request is converted into a feature vector. Third, the Random Forest is
fitted to the feature--label pairs. In detail:

\begin{enumerate}[leftmargin=1.8em, itemsep=2pt]
  \item Build the per-path parameter vocabulary from all training requests.
  \item Extract the $22$-feature vector for every labelled request, forming
        the feature matrix.
  \item Instantiate a Random Forest with $200$ trees and balanced class
        weights.
  \item Fit the Random Forest to the feature matrix and the binary labels
        ($1$ for attack, $0$ for benign).
  \item Persist the trained model together with the vocabulary.
\end{enumerate}

Recalling the mathematical formulation given in Chapter 2, the Random Forest
trains $B = 200$ decision trees, each grown on a bootstrap sample of the
training set and splitting nodes on a random subset of the $22$ features
using the Gini impurity $G = 1 - \sum_i p_i^2$ as the split criterion. The
key property exploited here is that the forest produces not only a hard label
but a calibrated probability, obtained by averaging the per-tree votes:
\begin{equation}
  p_{\text{attack}}(x) \;=\; \frac{1}{B}\sum_{b=1}^{B} \mathbf{1}\!\left[T_b(x) = \text{attack}\right],
  \label{eq:rfproba}
\end{equation}
where $T_b(x)$ is the prediction of the $b$-th tree. This is exactly the
quantity returned by the ``predict\_proba'' method and consumed by the
consensus rule. Figure~\ref{fig:code-rf} shows the training and inference
code of the supervised tier.

\begin{figure}[ht]
  \centering
  \fbox{\parbox{0.9\textwidth}{\centering\vspace{2.0cm}
    \textit{[Insert a screenshot of the ``train()'' and
    ``predict\_proba()'' methods of the ``SupervisedDetector''
    class from ``models/supervised\_detector.py'' here.]}\vspace{2.0cm}}}
  % Replace the box above with: \includegraphics[width=0.9\textwidth]{code_supervised.png}
  \caption{Source code of the supervised Random Forest tier
           (``supervised\_detector.py'').}
  \label{fig:code-rf}
\end{figure}

At inference time the model is queried for the probability
$p_{\text{attack}}$ that a request is malicious, computed by
Equation~\eqref{eq:rfproba}; the binary Tier~2 hit is
defined as $p_{\text{attack}} \geq 0.5$. Crucially, the consensus rule
(Section~\ref{sec:consensus}) consumes the continuous probability rather
than the thresholded decision, which lets it validate weaker anomaly signals
at a lower probability cut-off. On the evaluation set this tier attains a
precision of $94.80\%$, a recall of $94.76\%$, and an F1 score of $94.78\%$,
the highest of any single model in the system. It therefore acts as the
backbone of the hybrid: it both contributes a strong independent verdict and
serves as the arbiter that confirms or rejects the anomaly tier's findings.
Its inherent limitation, shared by all supervised methods, is that it can
only generalise from attack families present in the training data---hence
the need for the third tier.

% =====================================================================
\section{Tier 3 --- Unsupervised One-Class Anomaly Detection}
\label{sec:tier3}

The third tier addresses the blind spot of supervised learning by detecting
deviations from normal traffic without ever being shown an attack. As
explained in Chapter 2, it follows the one-class learning
principle: the model is trained \emph{exclusively on clean traffic}
($34{,}941$ benign log lines) so that it learns the distribution of normal
requests, and at inference any request that falls outside that distribution
is flagged as anomalous. Three algorithms were implemented---Local Outlier
Factor, Isolation Forest, and One-Class SVM---among which the Local Outlier
Factor was selected as the operational model owing to its superior
precision--recall balance.

Unlike the supervised tier, the training consumes no labels: the input is a
corpus of attack-free requests. The features are extracted with the shared
extractor, standardised so that distance-based density estimation is not
dominated by high-magnitude features, and then passed to the one-class
estimator. For the Local Outlier Factor the estimator is configured with
twenty neighbours and in novelty mode, which allows it to score requests
that were not part of the training corpus. The training steps are:

\begin{enumerate}[leftmargin=1.8em, itemsep=2pt]
  \item Build the per-path parameter vocabulary from the clean requests.
  \item Extract the $22$-feature vector for every clean request, forming the
        feature matrix.
  \item Fit a standardisation transform on the feature matrix and apply it
        so that all features share a comparable scale.
  \item Instantiate the Local Outlier Factor with twenty neighbours in
        novelty-detection mode.
  \item Fit the model on the standardised clean features so that it learns
        the ``normal'' manifold.
  \item Persist the model, the standardisation transform, and the
        vocabulary.
\end{enumerate}

The mathematical formulation of the three one-class estimators (the local
reachability density and Local Outlier Factor score, the isolation path
length and anomaly score, and the One-Class SVM margin objective) was given
in Chapter 2; Figure~\ref{fig:code-lof} shows how the chosen estimator is
trained on clean data in the ``train'' method of the
``LogAnomalyDetector'' class.

\begin{figure}[ht]
  \centering
  \fbox{\parbox{0.9\textwidth}{\centering\vspace{2.0cm}
    \textit{[Insert a screenshot of the ``train()'' method of the
    ``LogAnomalyDetector'' class (clean-only fitting of LOF/IF/OCSVM)
    from ``models/ai\_detector.py'' here.]}\vspace{2.0cm}}}
  % Replace the box above with: \includegraphics[width=0.9\textwidth]{code_oneclass.png}
  \caption{Source code of the unsupervised one-class training routine
           (``ai\_detector.py'').}
  \label{fig:code-lof}
\end{figure}

At inference the estimator returns an outlier indicator ($-1$ for an
outlier, $+1$ for an inlier); the tier reports an anomaly precisely when the
output is $-1$. The decisive role of the clean-only training discipline is
best illustrated by the experimental contrast: when the same model is
mistakenly trained on a $50/50$ mixture of attack and clean traffic, its F1
score collapses to about $33\%$, because the ``normal'' manifold becomes
polluted with attack points and the notion of an outlier loses meaning.
Trained correctly on clean data, the Local Outlier Factor reaches a
precision of $88.28\%$ and a recall of $95.13\%$ (F1 $=91.57\%$). The high
recall is exactly the property the hybrid needs from this tier: it catches
attacks the supervised model misses, at the cost of a higher false-positive
rate that the consensus rule is designed to suppress.

% =====================================================================
\section{Smart Hybrid Consensus and Threat Scoring}
\label{sec:consensus}

The final stage fuses the three tier outputs into a single decision. A naive
fusion would simply raise an alert whenever any tier fires (a logical OR);
however, because the Local Outlier Factor alone produces a non-negligible
number of false positives, such a rule would inflate the false-alarm rate.
The proposed \emph{smart consensus} instead lets the two high-precision
signals fire on their own, while treating the anomaly tier as a tripwire
that must be corroborated by the supervised model before it is trusted. Let
$R$ denote the set of Tier~1 rules that fired, $p$ the Random Forest attack
probability, and $a \in \{\text{true},\text{false}\}$ the Tier~3 anomaly
flag. A request is classified as an attack when
\begin{equation}
  \mathrm{attack} \;=\; (R \neq \emptyset) \;\lor\; (p \geq 0.5) \;\lor\;
  (a \,\land\, p \geq 0.3).
  \label{eq:consensus}
\end{equation}
The three disjuncts of Equation~\eqref{eq:consensus} encode three
independent ways to raise an alert. The first, $R \neq \emptyset$, fires when
any signature matched. The second, $p \geq 0.5$, fires when the Random Forest
is confident on its own. The third, $a \land p \geq 0.3$, is the key idea: a
Tier~3 anomaly is trusted only when the Random Forest is at least mildly
suspicious ($p \geq 0.3$), so an isolated anomaly that the Random Forest
flatly rejects ($p < 0.3$) is not alerted, which is what suppresses the
anomaly tier's false positives.

Independently of the binary verdict, the system assigns a continuous threat
score that drives the severity ranking shown to the analyst. The score
aggregates the three signals with weights reflecting their trustworthiness:
\begin{equation}
  \mathrm{score} \;=\; 0.45\,[\,R \neq \emptyset\,] \;+\;
  0.40\,p\,[\,p \geq 0.5\,] \;+\; 0.20\,[\,a\,],
  \label{eq:score}
\end{equation}
where the Iverson bracket $[\,\cdot\,]$ equals $1$ when its condition holds
and $0$ otherwise. The signature tier therefore contributes a fixed $0.45$
(the largest weight, as the most trusted signal), the supervised tier
contributes $0.40$ scaled by its own confidence $p$, and the anomaly tier
contributes a fixed $0.20$. The score is finally mapped to a severity level
by
\begin{equation}
  \mathrm{level} =
  \begin{cases}
    \textsc{critical} & \text{if } \mathrm{score} \geq 0.75,\\
    \textsc{high}     & \text{if } 0.45 \leq \mathrm{score} < 0.75,\\
    \textsc{medium}   & \text{if } 0.25 \leq \mathrm{score} < 0.45,\\
    \textsc{low}      & \text{if } \mathrm{score} < 0.25 \text{ and the verdict is an attack},\\
    \textsc{clean}    & \text{if the verdict is not an attack.}
  \end{cases}
  \label{eq:severity}
\end{equation}

Combining Equations~\eqref{eq:consensus} and~\eqref{eq:score} with the
severity mapping of Equation~\eqref{eq:severity}, the decision procedure
invoked for every request performs the following steps:

\begin{enumerate}[leftmargin=1.8em, itemsep=2pt]
  \item Initialise the threat score to zero.
  \item Add $0.45$ if any Tier~1 rule fired.
  \item Add $0.40 \times p$ if the Random Forest probability $p$ is at least
        $0.5$.
  \item Add $0.20$ if the Tier~3 anomaly flag is set.
  \item Compute the verdict using the smart consensus rule of
        Equation~\eqref{eq:consensus}.
  \item If the verdict is ``not an attack'', set the severity to
        \textsc{clean}; otherwise map the score to \textsc{critical},
        \textsc{high}, \textsc{medium}, or \textsc{low} using the thresholds
        of Equation~\eqref{eq:severity}.
  \item Return the verdict, the threat score, and the severity level.
\end{enumerate}

Figure~\ref{fig:code-consensus} shows the implementation of this fusion rule,
where the boolean verdict of Equation~\eqref{eq:consensus} is computed in a
single expression over the three tier outputs.

\begin{figure}[ht]
  \centering
  \fbox{\parbox{0.9\textwidth}{\centering\vspace{1.8cm}
    \textit{[Insert a screenshot of the smart-consensus decision logic
    (the ``should\_alert''/``is\_attack'' expression and the
    threat-score accumulation) from ``apache\_log.py'' here.]}\vspace{1.8cm}}}
  % Replace the box above with: \includegraphics[width=0.9\textwidth]{code_consensus.png}
  \caption{Source code of the smart hybrid consensus and threat-scoring logic
           in ``apache\_log.py''.}
  \label{fig:code-consensus}
\end{figure}

A practical advantage of this design is its \emph{explainability}: because
the verdict is a transparent combination of three named signals, every alert
can be annotated with the tiers that contributed to it (for example
``REGEX(SQL\_INJECTION) + SUPERVISED(RF, p=0.98)''), which helps the
analyst triage quickly. As reported for the full pipeline, the smart
consensus achieves a precision of $91.19\%$, a recall of $97.62\%$, an F1 of
$94.30\%$, and an F2 of $96.26\%$---the highest F2 in the system, confirming
that the fusion meets the IDS priority of minimising missed attacks while
keeping false alarms manageable.

% =====================================================================
\section{Operating Modes: Static Scan and Real-time Monitoring}
\label{sec:modes}

The same three-tier pipeline described above is exposed through two
operating modes, so that the system can be used both for the offline
analysis of historical logs and for live protection of a running web server.
Both modes share exactly the same parsing, feature extraction, three-tier
evaluation, and smart-consensus decision; they differ only in how log lines
enter the system and how the results are delivered.

\subsection{Static file scan}

The static-scan mode is intended for the offline, batch analysis of a log
file that has already been collected---for example for a forensic audit of
historical traffic or for the reproducible evaluation of the models. It is
invoked from the command line by passing the ``scan'' action, the path
to the log file, and the desired Tier-3 model, as in the command
\begin{quote}\small\ttfamily
python apache\_log.py scan <logfile> lof
\end{quote}
The system then reads the
file line by line and, for every request, runs the full pipeline and produces
two output files in JSON Lines format. The first, ``scan\_results.jsonl'',
contains one record \emph{per log line}---both clean and attack---each
carrying the verdict, the threat score, the severity level, and the list of
tiers that contributed, so that the outcome of every request can be audited.
The second, ``alerts.jsonl'', contains only the lines judged to be
attacks, in a compact alert format. At the end of the scan the system prints
a summary that reports the total number of lines processed, the split between
clean and attack, and the distribution of alerts across the
\textsc{critical}, \textsc{high}, \textsc{medium}, and \textsc{low} severity
levels. A representative record from ``scan\_results.jsonl'' is shown
below:

\begin{quote}\small\ttfamily
\{ "raw\_url": "/login?POST\_BODY=user...", "verdict": "ATTACK",\\
\ "level": "HIGH", "score": 0.94, "rule\_ids": [], \\
\ "sup\_prob": 0.98, "lof\_hit": true \}
\end{quote}

\subsection{Real-time monitoring}

The real-time monitoring mode turns the system into a live sensor that
analyses requests as they arrive. It reads log lines from standard input,
which lets it be attached directly to a growing server log with a single
shell pipeline, for example
\begin{quote}\footnotesize\ttfamily
tail -F /var/log/apache2/access.log |\\
\hspace*{1.5em}python apache\_log.py monitor /dev/null lof
\end{quote}
Every incoming line is passed through the identical
three-tier consensus, and whenever a request is judged to be an attack the
system emits a colour-coded alert on the console---the colour reflecting the
severity level---showing the timestamp, the source IP address, the aggregated
threat score, the contributing tiers, and the offending request. In parallel,
each alert is appended as a structured record to
``runtime/monitor\_alerts.jsonl''. A companion dashboard, implemented
with Streamlit, continuously tails this file and presents the information
visually: key-performance-indicator cards (such as the detection rate and the
false-positive rate), a live feed of the most recent alerts, the most active
source IP addresses, and the distribution of alerts by threat level. Together,
the console output and the dashboard give an analyst an immediate, continuously
updated picture of the attacks hitting the server.
Figure~\ref{fig:alert-monitor} shows the console of the real-time monitor at
the moment an attack is detected, and Figure~\ref{fig:dashboard} shows the
corresponding view on the dashboard.

\begin{figure}[ht]
  \centering
  \fbox{\parbox{0.9\textwidth}{\centering\vspace{2.2cm}
    \textit{[Insert a screenshot of the real-time monitor console at the
    moment an alert is raised (the colour-coded ``[CRITICAL]/[HIGH]''
    line with IP, threat score, tiers, and request) here.]}\vspace{2.2cm}}}
  % Replace the box above with: \includegraphics[width=0.9\textwidth]{screenshot_monitor_alert.png}
  \caption{Real-time monitoring console raising an alert the moment an attack
           is detected.}
  \label{fig:alert-monitor}
\end{figure}

\begin{figure}[ht]
  \centering
  \fbox{\parbox{0.9\textwidth}{\centering\vspace{2.2cm}
    \textit{[Insert a screenshot of the Streamlit dashboard (KPI cards, live
    alert feed, top source IPs, threat-level distribution) here.]}\vspace{2.2cm}}}
  % Replace the box above with: \includegraphics[width=0.9\textwidth]{screenshot_dashboard.png}
  \caption{Real-time monitoring dashboard summarising the live alert feed and
           key performance indicators.}
  \label{fig:dashboard}
\end{figure}

% =====================================================================
\section{Chapter Summary}
\label{sec:ch3summary}

This chapter has presented the proposed solution in full. Section~\ref{sec:overview}
gave an overview of the three-tier hybrid pipeline and described its data
flow (Figure~\ref{fig:pipeline}) together with the overall scanning
procedure. The subsequent sections detailed each stage: request parsing with
multi-layer decoding and the construction of the shared $22$-feature vector
(Section~\ref{sec:pipeline-feat}); the always-on signature and rule tier with
its defensive decoding and stateful tampering rules (Section~\ref{sec:tier1});
the supervised Random Forest tier that supplies a calibrated attack
probability (Section~\ref{sec:tier2}); and the unsupervised one-class anomaly
tier trained exclusively on clean traffic (Section~\ref{sec:tier3}). Finally,
Section~\ref{sec:consensus} introduced the smart consensus rule
(Equation~\eqref{eq:consensus}) and the weighted threat score
(Equation~\eqref{eq:score}), which together turn three complementary opinions
into a single, explainable, and recall-oriented verdict. Finally,
Section~\ref{sec:modes} described the two operating modes that expose this
pipeline in practice: the static file scan for offline, batch analysis of
historical logs, and the real-time monitor with its live dashboard for
protecting a running server. The next chapter evaluates this solution
empirically, comparing the individual tiers and the hybrid configurations on
the benchmark dataset.

% --------------------- STANDALONE WRAPPER (begin) --------------------
\end{document}
% ---------------------- STANDALONE WRAPPER (end) ---------------------
